\FigureLayoutDeclare{img/2006/tianjin_science/q19_fig1.png}{8.0cm}{29bp 21bp 59bp 24bp}
\FigureLayoutDeclare{img/2006/tianjin_science/q22_fig1.png}{9.0cm}{90bp 7bp 61bp 20bp}

\chapter{2006年天津卷（理）}
\section{单选题}
\begin{problem}
\(\i\) 是虚数单位，\(\frac{\i}{1+\i}=\)（\quad）

\choices
    {\(\frac{1}{2}+\frac{1}{2}\i\)}
    {\(-\frac{1}{2}+\frac{1}{2}\i\)}
    {\(\frac{1}{2}-\frac{1}{2}\i\)}
    {\(-\frac{1}{2}-\frac{1}{2}\i\)}
\end{problem}

\begin{answer}
A
\end{answer}

\begin{solution}
有理化分母得
\[
\frac{\i}{1+\i}=\frac{\i(1-\i)}{(1+\i)(1-\i)}=\frac{1+\i}{2}=\frac{1}{2}+\frac{1}{2}\i.
\]
故选 A.
\end{solution}

\begin{problem}
如果双曲线的两个焦点分别为 \(F_{1}(-3,0)\)，\(F_{2}(3,0)\)，一条渐近线方程为 \(y=\sqrt{2}x\)，那么它的两条渐近线间的距离是（\quad）

\choices
    {\(6\sqrt{3}\)}
    {\(4\)}
    {\(2\)}
    {\(1\)}
\end{problem}

\begin{answer}
C
\end{answer}

\begin{solution}
原试卷题干写作“两条渐近线间的距离”，但双曲线的两条渐近线相交，按字面不能得到选项中的非零值；结合参考答案中的第三项，以下按题意为“两条准线间的距离”求解. 设双曲线的标准方程为 \(\dfrac{x^{2}}{a^{2}}-\dfrac{y^{2}}{b^{2}}=1\)，焦点在 \(x\) 轴上，则其渐近线为 \(y=\pm\dfrac{b}{a}x\). 由题意，\(\dfrac{b}{a}=\sqrt{2}\)，又焦半距为 \(c=3\)，所以 \(b^{2}=2a^{2}\)，\(a^{2}+b^{2}=c^{2}=9\).
从而 \(a^{2}=3\)，\(b^{2}=6\). 双曲线的离心率为 \(e=\dfrac{c}{a}\)，两条准线为 \(x=\dfrac{a}{e}=\dfrac{a^{2}}{c}=1\) 和 \(x=-\dfrac{a}{e}=-\dfrac{a^{2}}{c}=-1\)，故两条准线间的距离为 \(2\). 因此选 C.
\end{solution}

\begin{problem}
设变量 \(x\)，\(y\) 满足约束条件 \(\begin{cases}
y\leqslant x,\\
x+y\geqslant 2,\\
y\geqslant 3x-6,
\end{cases}\) 则目标函数 \(z=2x+y\) 的最小值为（\quad）

\choices
    {\(2\)}
    {\(3\)}
    {\(4\)}
    {\(9\)}
\end{problem}

\begin{answer}
B
\end{answer}

\begin{solution}
由 \(x+y\geqslant2\) 与 \(y\leqslant x\) 得 \(x\geqslant1\)，由 \(y\geqslant3x-6\) 与 \(y\leqslant x\) 得 \(x\leqslant3\)，故可行域为有界多边形. 其顶点由三条边界直线 \(y=x\)，\(x+y=2\)，\(y=3x-6\) 两两相交得到，分别为 \((1,1)\)，\((2,0)\)，\((3,3)\)，且这三个点都满足约束条件. 目标函数 \(z=2x+y\) 在三点处的值分别为 \(3\)，\(4\)，\(9\)，所以最小值为 \(3\). 故选 B.
\end{solution}

\begin{problem}
设集合 \(M=\{x\mid 0<x\leqslant 3\}\)，\(N=\{x\mid 0<x\leqslant 2\}\)，那么“\(a\in M\)”是“\(a\in \N\)”的（\quad）

\choices
    {充分而不必要条件}
    {必要而不充分条件}
    {充分必要条件}
    {既不充分也不必要条件}
\end{problem}

\begin{answer}
B
\end{answer}

\begin{solution}
因为 \(N=(0,2]\subset M=(0,3]\)，所以 \(a\in N\) 能推出 \(a\in M\)，但 \(a\in M\) 不能推出 \(a\in N\)，例如 \(a=\dfrac{5}{2}\) 时结论不成立. 因而“\(a\in M\)”是“\(a\in N\)”的必要而不充分条件，故选 B.
\end{solution}

\begin{problem}
将 \(4\) 个颜色互不相同的球全部放入编号为 \(1\) 和 \(2\) 的两个盒子里，使得放入每个盒子里的球的个数不小于该盒子的编号，则不同的放球方法有（\quad）

\choices
    {\(10\) 种}
    {\(20\) 种}
    {\(36\) 种}
    {\(52\) 种}
\end{problem}

\begin{answer}
A
\end{answer}

\begin{solution}
设放入编号为 \(2\) 的盒子的球数为 \(k\). 由题意，\(k\geqslant2\)，且另一个盒子至少有 \(1\) 个球，所以 \(k\) 只能取 \(2\) 或 \(3\). 于是放球方法数为
\[
\mathrm C_{4}^{2}+\mathrm C_{4}^{3}=6+4=10.
\]
故选 A.
\end{solution}

\begin{problem}
设 \(m\)，\(n\) 是两条不同的直线，\(\alpha\)，\(\beta\) 是两个不同的平面. 考查下列命题，其中正确的命题是（\quad）

\choices
    {\(m\perp\alpha, n\subset\beta, m\perp n \Rightarrow \alpha\perp\beta\)}
    {\(\alpha\parallel\beta\)，\(m\perp\alpha\)，\(n\parallel\beta \Rightarrow m\perp n\)}
    {\(\alpha\perp\beta\)，\(m\perp\alpha\)，\(n\parallel\beta \Rightarrow m\perp n\)}
    {\(\alpha\perp\beta\)，\(\alpha\cap\beta=m\)，\(n\perp m \Rightarrow n\perp\beta\)}
\end{problem}

\begin{answer}
B
\end{answer}

\begin{solution}
若 \(\alpha\parallel\beta\)，直线 \(m\perp\alpha\)，则 \(m\) 也垂直于平面 \(\beta\). 又 \(n\parallel\beta\)，所以 \(n\) 的方向与平面 \(\beta\) 内某条直线的方向相同，从而 \(m\perp n\). 因此命题 B 正确，故选 B.
\end{solution}

\begin{problem}
已知数列 \(\{a_{n}\}\)，\(\{b_{n}\}\) 都是公差为 \(1\) 的等差数列，其首项分别为 \(a_{1}\)，\(b_{1}\)，且 \(a_{1}+b_{1}=5\)，\(a_{1}\)，\(b_{1}\in\N^{*}\). 设 \(C_{n}=a_{b_{n}}\)（\(n\in\N^{*}\)），则数列 \(\{C_{n}\}\) 的前 \(10\) 项和等于（\quad）

\choices
    {\(55\)}
    {\(70\)}
    {\(85\)}
    {\(100\)}
\end{problem}

\begin{answer}
C
\end{answer}

\begin{solution}
由两数列的公差都为 \(1\)，有
\(a_{k}=a_{1}+k-1\)，\(b_{n}=b_{1}+n-1\).
因此
\[
C_{n}=a_{b_{n}}=a_{1}+b_{n}-1=a_{1}+b_{1}+n-2=n+3.
\]
所以前 \(10\) 项和为
\[
\sum_{n=1}^{10}C_{n}=\sum_{n=1}^{10}(n+3)=\frac{(4+13)\cdot10}{2}=85.
\]
故选 C.
\end{solution}

\begin{problem}
已知函数 \(f(x)=a\sin x-b\cos x\)（\(a\)，\(b\) 为常数，\(a\neq 0\)，\(x\in\R\)）在 \(x=\frac{\pi}{4}\) 处取得最小值，则函数 \(y=f\left(\frac{3\pi}{4}-x\right)\) 是（\quad）

\choices
    {偶函数且它的图象关于点 \((\pi,0)\) 对称}
    {偶函数且它的图象关于点 \(\left(\frac{3\pi}{2},0\right)\) 对称}
    {奇函数且它的图象关于点 \(\left(\frac{3\pi}{2},0\right)\) 对称}
    {奇函数且它的图象关于点 \((\pi,0)\) 对称}
\end{problem}

\begin{answer}
D
\end{answer}

\begin{solution}
设 \(R=\sqrt{a^{2}+b^{2}}>0\). 由于 \(f(x)\) 在 \(x=\dfrac{\pi}{4}\) 处取得最小值，可将它写成
\[
f\left(\frac{\pi}{4}+t\right)=-R\cos t.
\]
令 \(h(x)=f\left(\dfrac{3\pi}{4}-x\right)\)，则
\[
h(x)=-R\cos\left(\frac{\pi}{2}-x\right)=-R\sin x.
\]
所以 \(h(-x)=-h(x)\)，它是奇函数；又 \(h(2\pi-x)=-h(x)\)，故其图象关于点 \((\pi,0)\) 对称. 因此选 D.
\end{solution}

\begin{problem}
函数 \(f(x)\) 的定义域为开区间 \((a,b)\)，导函数 \(f'(x)\) 在 \((a,b)\) 内的图像如图所示，则函数 \(f(x)\) 在开区间 \((a,b)\) 内有极小值点（\quad）

\bitmapfigure[max width=0.60\linewidth]{img/2006/tianjin_science/q9_fig1.png}

\choices
    {\(1\) 个}
    {\(2\) 个}
    {\(3\) 个}
    {\(4\) 个}
\end{problem}

\begin{answer}
A
\end{answer}

\begin{solution}
从图象看，\(f'(x)\) 在定义域内共有四处零点：除原点外的三处零点分别由正变负、由负变正、由正变负；原点处虽为零点，但两侧均为正，符号没有改变. 函数取得极小值的条件是导数由负变正，因此只有第二处零点对应极小值点，共 \(1\) 个. 故选 A.
\end{solution}

\begin{problem}
已知函数 \(y=f(x)\) 的图象与函数 \(y=a^{x}\)（\(a>0\) 且 \(a\neq 1\)）的图象关于直线 \(y=x\) 对称，记 \(g(x)=f(x)[f(x)+f(2)-1]\). 若 \(y=g(x)\) 在区间 \(\left[\frac{1}{2},2\right]\) 上是增函数，则实数 \(a\) 的取值范围是（\quad）

\choices
    {\(\left[2,+\infty\right)\)}
    {\((0,1)\cup(1,2)\)}
    {\(\left[\frac{1}{2},1\right)\)}
    {\(\left(0,\frac{1}{2}\right]\)}
\end{problem}

\begin{answer}
D
\end{answer}

\begin{solution}
由反函数关系，\(f(x)=\log_{a}x\). 令 \(q=\log_{a}2\)，则
\[
g'(x)=\frac{1}{x\ln a}\left(2\log_{a}x+\log_{a}2-1\right).
\]
当 \(a>1\) 时，\(\ln a>0\)，而在 \(x=\dfrac12\) 处括号内的值为 \(-q-1<0\)，故在区间左端附近 \(g'(x)<0\)，不可能为增函数.

当 \(0<a<1\) 时，\(\ln a<0\)，所以要使 \(g\) 在 \(\left[\dfrac12,2\right]\) 上递增，需使
\[
2\log_{a}x+q-1\leqslant0.
\]
因为 \(\log_{a}x\) 在该区间上递减，括号内的最大值在 \(x=\dfrac12\) 处取得，故条件等价于 \(-q-1\leqslant0\)，即 \(q\geqslant-1\). 由 \(q=\dfrac{\ln2}{\ln a}\) 且 \(0<a<1\)，得到 \(0<a\leqslant\dfrac12\). 因此选 D.
\end{solution}

\section{填空题}
\begin{problem}
\(\left(2x+\frac{1}{\sqrt{x}}\right)^{7}\) 的二项展开式中 \(x\) 的系数是\fillinblank{}. （用数字作答）
\end{problem}

\begin{answer}
\(280\)
\end{answer}

\begin{solution}
展开式的通项为
\[
\mathrm C_{7}^{k}(2x)^{k}\left(x^{-\frac12}\right)^{7-k}=\mathrm C_{7}^{k}2^{k}x^{\frac{3k-7}{2}}.
\]
令 \(\dfrac{3k-7}{2}=1\)，得 \(k=3\)，所以 \(x\) 的系数为
\[
\mathrm C_{7}^{3}2^{3}=35\cdot8=280.
\]
\end{solution}

\begin{problem}
设向量 \(\bs{a}\) 与 \(\bs{b}\) 的夹角为 \(\theta\)，且 \(\bs{a}=(3,3)\)，\(2\bs{b}-\bs{a}=(-1,1)\)，则 \(\cos\theta=\)\fillinblank{}.
\end{problem}

\begin{answer}
\(\frac{3\sqrt{10}}{10}\)
\end{answer}

\begin{solution}
由 \(2\bs{b}-\bs{a}=(-1,1)\) 和 \(\bs{a}=(3,3)\)，得
\[
2\bs{b}=(2,4).
\]
所以 \(\bs{b}=(1,2)\). 于是
\[
\bs{a}\cdot\bs{b}=3\cdot1+3\cdot2=9.
\]
又 \(|\bs{a}|=3\sqrt2\)，\(|\bs{b}|=\sqrt5\).
因此
\[
\cos\theta=\frac{\bs{a}\cdot\bs{b}}{|\bs{a}|\,|\bs{b}|}=\frac{9}{3\sqrt{10}}=\frac{3\sqrt{10}}{10}.
\]
\end{solution}

\begin{problem}
如图，在正三棱柱 \(ABC-A_{1}B_{1}C_{1}\) 中，\(AB=1\). 若二面角 \(C-AB-C_{1}\) 的大小为 \(60^{\circ}\)，则点 \(C\) 到平面 \(ABC_{1}\) 的距离为\fillinblank{}.

\bitmapfigure[float=inline,max width=0.42\linewidth]{img/2006/tianjin_science/q13_fig1.png}
\FloatBarrier
\end{problem}

\begin{answer}
\(\frac{3}{4}\)
\end{answer}

\begin{solution}
建立直角坐标系，取 \(AB\) 为 \(x\) 轴，令底面 \(ABC\) 在 \(xOy\) 平面内，设棱柱的高为 \(h\)，则可取
\(A=(0,0,0)\)，\(B=(1,0,0)\)，\(C=\left(\frac12,\frac{\sqrt3}{2},0\right)\)，\(C_{1}=\left(\frac12,\frac{\sqrt3}{2},h\right)\).
设 \(M\) 为 \(AB\) 的中点. 在过 \(M\) 且垂直于 \(AB\) 的截面内，二面角 \(C-AB-C_{1}\) 就是直线 \(CM\) 与直线 \(C_{1}M\) 的夹角，所以
\[
\cos60^{\circ}=\frac{\frac{\sqrt3}{2}}{\sqrt{\frac34+h^{2}}}.
\]
由此得 \(h=\dfrac32\). 平面 \(ABC_{1}\) 的方程为
\[
3y-\sqrt3z=0.
\]
所以点 \(C\) 到该平面的距离为
\[
\frac{\left|3\cdot\frac{\sqrt3}{2}-\sqrt3\cdot0\right|}{\sqrt{3^{2}+\left(-\sqrt3\right)^{2}}}=\frac34.
\]
\end{solution}

\begin{problem}
设直线 \(ax-y+3=0\) 与圆 \((x-1)^2+(y-2)^2=4\) 相交于 \(A\)，\(B\) 两点，且弦 \(AB\) 的长为 \(2\sqrt{3}\)，则 \(a=\)\fillinblank{}.
\end{problem}

\begin{answer}
\(0\)
\end{answer}

\begin{solution}
圆心为 \(O(1,2)\)，半径为 \(2\). 弦长为 \(2\sqrt3\)，所以圆心到弦所在直线的距离为
\[
d=\sqrt{2^{2}-\left(\frac{2\sqrt3}{2}\right)^{2}}=1.
\]
直线 \(ax-y+3=0\) 到圆心的距离也应为 \(1\)，故
\[
\frac{|a-2+3|}{\sqrt{a^{2}+1}}=1.
\]
两边平方得 \((a+1)^{2}=a^{2}+1\)，从而 \(a=0\).
\end{solution}

\begin{problem}
某公司一年购买某种货物 \(400\text{ 吨}\)，每次都购买 \(x\text{ 吨}\)，运费为 \(4\text{ 万元/次}\)，一年的总存储费用为 \(4x\text{ 万元}\)，要使一年的总运费与总存储费用之和最小，则 \(x=\)\fillinblank{}\(\text{吨}\).
\end{problem}

\begin{answer}
\(20\)
\end{answer}

\begin{solution}
一年购买次数为 \(\dfrac{400}{x}\)，所以一年总运费与总存储费用之和为
\[
T(x)=4\cdot\frac{400}{x}+4x=\frac{1600}{x}+4x,\qquad x>0.
\]
由基本不等式，
\[
T(x)=\frac{1600}{x}+4x\geqslant2\sqrt{\frac{1600}{x}\cdot4x}=160,
\]
等号成立当且仅当 \(\dfrac{1600}{x}=4x\)，即 \(x=20\). 因此总费用最小时 \(x=20\) 吨.
\end{solution}

\begin{problem}
设函数 \(f(x)=\frac{1}{x+1}\)，点 \(A_{0}\) 表示坐标原点，点 \(A_{n}(n,f(n))\)（\(n\in\N^{*}\)）. 若向量 \(\bs{a}_{n}=\overrightarrow{A_{0}A_{1}}+\overrightarrow{A_{1}A_{2}}+\cdots+\overrightarrow{A_{n-1}A_{n}}\)，\(\theta_{n}\) 是 \(\bs{a}_{n}\) 与 \(\bs{i}\) 的夹角（其中 \(\bs{i}=(1,0)\)），设 \(S_{n}=\tan\theta_{1}+\tan\theta_{2}+\cdots+\tan\theta_{n}\)，则 \(\displaystyle\lim_{n\to\infty}S_{n}=\)\fillinblank{}.
\end{problem}

\begin{answer}
\(1\)
\end{answer}

\begin{solution}
向量首尾相接后中间项相消，故
\[
\bs{a}_{n}=\overrightarrow{A_{0}A_{n}}=\left(n,\frac{1}{n+1}\right).
\]
由于 \(\bs{i}=(1,0)\)，且 \(\bs{a}_{n}\) 的横坐标为正，有
\[
\tan\theta_{n}=\frac{\frac{1}{n+1}}{n}=\frac{1}{n(n+1)}=\frac1n-\frac1{n+1}.
\]
于是
\[
S_{n}=\sum_{k=1}^{n}\left(\frac1k-\frac1{k+1}\right)=1-\frac1{n+1},
\]
从而 \(\displaystyle\lim_{n\to\infty}S_{n}=1\).
\end{solution}

\section{解答题}
\begin{problem}
如图，在 \(\triangle ABC\) 中，\(AC=2\)，\(BC=1\)，\(\cos C=\frac{3}{4}\).

\begin{enumerate}
\item 求 \(AB\) 的值；
\item 求 \(\sin(2A+C)\) 的值.
\end{enumerate}

\bitmapfigure[max width=0.34\linewidth]{img/2006/tianjin_science/q17_fig1.png}
\end{problem}

\begin{solution}
\begin{enumerate}
\item 由余弦定理，
\[
AB^{2}=AC^{2}+BC^{2}-2AC\cdot BC\cos C=2^{2}+1^{2}-2\cdot2\cdot1\cdot\frac34=2,
\]
所以 \(AB=\sqrt2\).
\item 由三角形内角和，\(2A+C=\pi+A-B\)，从而
\[
\sin(2A+C)=-\sin(A-B)=\sin(B-A).
\]
再次由余弦定理，
\[
\cos A=\frac{AB^{2}+AC^{2}-BC^{2}}{2AB\cdot AC}=\frac{5\sqrt2}{8}.
\]
同理，\(\cos B=\frac{AB^{2}+BC^{2}-AC^{2}}{2AB\cdot BC}=-\frac{\sqrt2}{4}\).
由于 \(A\)，\(B\) 是三角形内角，故
\[
\sin A=\sqrt{1-\cos^{2}A}=\frac{\sqrt{14}}{8}.
\]
同理，\(\sin B=\sqrt{1-\cos^{2}B}=\frac{\sqrt{14}}{4}\).
因此
\[
\sin(B-A)=\sin B\cos A-\cos B\sin A=\frac{3\sqrt7}{8}.
\]
所以 \(\sin(2A+C)=\dfrac{3\sqrt7}{8}\).
\end{enumerate}
\end{solution}

\begin{problem}
某射手进行射击训练，假设每次射击击中目标的概率为 \(\frac{3}{5}\)，且每次射击的结果互不影响.

\begin{enumerate}
\item 求射手在 \(3\) 次射击中，至少有两次连续击中目标的概率（用数字作答）；
\item 求射手第 \(3\) 次击中目标时，恰好射击了 \(4\) 次的概率（用数字作答）；
\item 设随机变量 \(\xi\) 表示射手第 \(3\) 次击中目标时已射击的次数，求 \(\xi\) 的分布列.
\end{enumerate}
\end{problem}

\begin{solution}
记击中目标为 \(H\)，未击中目标为 \(M\)，则 \(P(H)=\dfrac35\)，\(P(M)=\dfrac25\).
\begin{enumerate}
\item 在三次射击中至少有两次连续击中，可能的结果为 \(HHH\)，\(HHM\)，\(MHH\)，故所求概率为
\[
P(HHH)+P(HHM)+P(MHH)=\left(\frac35\right)^{3}+2\left(\frac35\right)^{2}\frac25=\frac{63}{125}.
\]
\item 第四次射击恰好是第 \(3\) 次击中，等价于前 \(3\) 次中恰有 \(2\) 次击中且第 \(4\) 次击中，因此概率为
\[
\mathrm C_{3}^{2}\left(\frac35\right)^{2}\frac25\cdot\frac35=\frac{162}{625}.
\]
\item 当 \(k=3,4,\cdots\) 时，事件 \(\{\xi=k\}\) 表示前 \(k-1\) 次中恰有 \(2\) 次击中，第 \(k\) 次击中，所以
\[
P(\xi=k)=\mathrm C_{k-1}^{2}\left(\frac35\right)^{3}\left(\frac25\right)^{k-3}.
\]
因此 \(\xi\) 的分布列可写为
\begin{center}
\begin{tabular}{c|c}
\hline
\(\xi\) & \(3,4,5,\cdots\) \\
\hline
\(P\) & \(\mathrm C_{k-1}^{2}\left(\dfrac35\right)^{3}\left(\dfrac25\right)^{k-3}\ (k=3,4,5,\cdots)\) \\
\hline
\end{tabular}
\end{center}
\end{enumerate}
\end{solution}

\begin{problem}
如图，在五面体 \(ABCDEF\) 中，点 \(O\) 是矩形 \(ABCD\) 的对角线的交点，面 \(CDE\) 是等边三角形，棱 \(EF\parallel BC\) 且 \(EF=\frac{1}{2}BC\).

\begin{enumerate}
\item 证明：\(FO\parallel\) 平面 \(CDE\)；
\item 设 \(BC=\sqrt{3}CD\)，证明：\(EO\perp\) 平面 \(CDF\).
\end{enumerate}

\bitmapfigure[max width=0.55\linewidth]{img/2006/tianjin_science/q19_fig1.png}
\end{problem}

\begin{solution}
取 \(D\) 为原点，记
\(\bs{u}=\overrightarrow{DC}=\overrightarrow{AB}\)，\(\bs{v}=\overrightarrow{BC}\)，\(\bs{w}=\overrightarrow{DE}\).
因为 \(ABCD\) 是矩形，所以 \(\bs{u}\perp\bs{v}\). 又因为 \(CDE\) 是等边三角形，所以
\(|\bs{w}|=|\bs{u}|\)，\(|\bs{w}-\bs{u}|=|\bs{u}|\)，
从而 \(\bs{u}\cdot\bs{w}=\dfrac12|\bs{u}|^{2}\). 由图中棱的方向，\(\overrightarrow{FE}=\dfrac12\bs{v}\)，于是
\(\overrightarrow{DF}=\bs{w}-\frac12\bs{v}\)，\(\overrightarrow{DO}=\frac12(\bs{u}-\bs{v})\).
\begin{enumerate}
\item 有
\[
\overrightarrow{FO}=\overrightarrow{DO}-\overrightarrow{DF}=\frac12\bs{u}-\bs{w}.
\]
因为 \(\bs{u}\) 和 \(\bs{w}\) 都是平面 \(CDE\) 内的方向向量，所以 \(FO\parallel\) 平面 \(CDE\).
\item 由 \(BC=\sqrt3CD\)，得 \(|\bs{v}|^{2}=3|\bs{u}|^{2}\). 又
\[
\overrightarrow{OE}=\bs{w}-\frac12(\bs{u}-\bs{v})=\bs{w}-\frac12\bs{u}+\frac12\bs{v}.
\]
先有
\[
\overrightarrow{OE}\cdot\bs{u}=\bs{w}\cdot\bs{u}-\frac12|\bs{u}|^{2}+\frac12\bs{v}\cdot\bs{u}=0,
\]
所以 \(OE\perp CD\). 再计算
\[
\begin{aligned}
\overrightarrow{OE}\cdot\overrightarrow{DF}
&=\left(\bs{w}-\frac12\bs{u}+\frac12\bs{v}\right)\cdot\left(\bs{w}-\frac12\bs{v}\right)\\
&=|\bs{w}|^{2}-\frac12\bs{u}\cdot\bs{w}-\frac14|\bs{v}|^{2}\\
&=|\bs{u}|^{2}-\frac14|\bs{u}|^{2}-\frac14|\bs{v}|^{2}=0.
\end{aligned}
\]
所以 \(OE\perp DF\). 直线 \(CD\) 与 \(DF\) 在平面 \(CDF\) 内相交，故 \(EO\perp\) 平面 \(CDF\).
\end{enumerate}
\end{solution}

\begin{problem}
已知函数 \(f(x)=4x^{3}-3x^{2}\cos\theta+\frac{3}{16}\cos\theta\)，其中 \(x\in\R\)，\(\theta\) 为参数，且 \(0\leqslant\theta<2\pi\).

\begin{enumerate}
\item 当 \(\cos\theta=0\) 时，判断函数 \(f(x)\) 是否有极值；
\item 要使函数 \(f(x)\) 的极小值大于零，求参数 \(\theta\) 的取值范围；
\item 若对（2）中所求的取值范围内的任意参数 \(\theta\)，函数 \(f(x)\) 在区间 \((2a-1,a)\) 内都是增函数，求实数 \(a\) 的取值范围.
\end{enumerate}
\end{problem}

\begin{solution}
记 \(c=\cos\theta\)，则
\[
f'(x)=12x^{2}-6cx=6x(2x-c).
\]
\begin{enumerate}
\item 当 \(c=0\) 时，\(f(x)=4x^{3}\)，且 \(f'(x)=12x^{2}\geqslant0\). 函数在整个实数范围内单调递增，虽然 \(f'(0)=0\)，但导数没有由正变负或由负变正，故函数没有极值.
\item 当 \(c<0\) 时，\(x=0\) 是极小值点，但 \(f(0)=\dfrac{3c}{16}<0\)，不符合题意. 当 \(c>0\) 时，\(x=0\) 是极大值点，\(x=\dfrac{c}{2}\) 是极小值点，且
\[
f\left(\frac{c}{2}\right)=4\left(\frac{c}{2}\right)^{3}-3\left(\frac{c}{2}\right)^{2}c+\frac{3c}{16}=\frac{c(3-4c^{2})}{16}.
\]
要使极小值大于零，需且只需 \(0<c<\dfrac{\sqrt3}{2}\). 结合 \(0\leqslant\theta<2\pi\)，得
\[
\theta\in\left(\frac{\pi}{6},\frac{\pi}{2}\right)\cup\left(\frac{3\pi}{2},\frac{11\pi}{6}\right).
\]
\item 对上述范围内的 \(\theta\)，有 \(0<c<\dfrac{\sqrt3}{2}\)，所以 \(f'(x)>0\) 的区间为 \(\left(-\infty,0\right)\) 和 \(\left(\dfrac{c}{2},+\infty\right)\). 要使 \(f(x)\) 在 \((2a-1,a)\) 内对所有这些参数都递增，要么 \(a\leqslant0\)，要么对所有 \(0<c<\dfrac{\sqrt3}{2}\) 都有 \(2a-1\geqslant\dfrac{c}{2}\).
第二种情形等价于 \(2a-1\geqslant\dfrac{\sqrt3}{4}\)，并且区间非空要求 \(a<1\). 因此
\[
a\in\left(-\infty,0\right]\cup\left[\frac12+\frac{\sqrt3}{8},1\right).
\]
\end{enumerate}
\end{solution}

\begin{problem}
已知数列 \(\{x_{n}\}\)，\(\{y_{n}\}\) 满足 \(x_{1}=x_{2}=1\)，\(y_{1}=y_{2}=2\)，并且 \(\frac{x_{n+1}}{x_n}=\lambda\frac{x_n}{x_{n-1}}\)，\(\frac{y_{n+1}}{y_n}\geqslant\lambda\frac{y_n}{y_{n-1}}\)（\(\lambda\) 为非零参数，\(n=2\)，\(3\)，\(4\)，\(\cdots\)）.

\begin{enumerate}
\item 若 \(x_{1}\)，\(x_{3}\)，\(x_{5}\) 成等比数列，求参数 \(\lambda\) 的值；
\item 当 \(\lambda>0\) 时，证明：\(\frac{x_{n+1}}{y_{n+1}}\leqslant\frac{x_n}{y_n}\)（\(n\in\N^{*}\)）；
\item 当 \(\lambda>1\)，证明：\(\frac{x_{1}-y_{1}}{x_{2}-y_{2}}+\frac{x_{2}-y_{2}}{x_{3}-y_{3}}+\cdots+\frac{x_{n}-y_{n}}{x_{n+1}-y_{n+1}}<\frac{\lambda}{\lambda-1}\)（\(n\in\N^{*}\)）.
\end{enumerate}
\end{problem}

\begin{solution}
\begin{enumerate}
\item 由递推关系逐项计算得 \(x_{3}=\lambda\)，\(x_{4}=\lambda^{3}\)，\(x_{5}=\lambda^{6}\). 因为 \(x_{1}\)，\(x_{3}\)，\(x_{5}\) 成等比数列，所以 \(x_{3}^{2}=x_{1}x_{5}\)，即 \(\lambda^{2}=\lambda^{6}\).
又 \(\lambda\neq0\)，故 \(\lambda=\pm1\).
\item 当 \(n=1\) 时，\(\dfrac{x_{2}}{y_{2}}=\dfrac{x_{1}}{y_{1}}=\dfrac12\). 对 \(n\geqslant2\)，令 \(r_{n}=\dfrac{x_{n}}{x_{n-1}}\)，\(s_{n}=\dfrac{y_{n}}{y_{n-1}}\). 由 \(r_{2}=1\) 及递推式，
得 \(r_{n}=\lambda^{n-2}\)，\(\dfrac{x_{n+1}}{x_{n}}=\lambda^{n-1}\).
当 \(\lambda>0\) 时，由 \(s_{2}=1\) 和 \(s_{n+1}\geqslant\lambda s_{n}\) 归纳可得 \(s_{n+1}\geqslant\lambda^{n-1}\). 因为各项均为正数，所以
\[
\frac{x_{n+1}}{y_{n+1}}=\frac{x_{n}}{y_{n}}\cdot\frac{x_{n+1}/x_{n}}{y_{n+1}/y_{n}}\leqslant\frac{x_{n}}{y_{n}}.
\]
\item 当 \(\lambda>1\) 时，由上面的估计和初值可知 \(y_{n}>x_{n}>0\). 令 \(d_{n}=y_{n}-x_{n}>0\)，则 \(d_{1}=d_{2}=1\)，且对 \(n\geqslant2\)，
有 \(y_{n+1}\geqslant\lambda^{n-1}y_{n}\)，\(x_{n+1}=\lambda^{n-1}x_{n}\)，
从而
\[
d_{n+1}\geqslant\lambda^{n-1}d_{n}.
\]
因此
\[
\frac{x_{n}-y_{n}}{x_{n+1}-y_{n+1}}=\frac{d_{n}}{d_{n+1}}\leqslant\frac{1}{\lambda^{n-1}}\qquad(n\geqslant2).
\]
当 \(n=1\) 时，该比值为 \(d_{1}/d_{2}=1\)，同样不超过 \(1/\lambda^{0}\). 因此相加得到
\[
\sum_{k=1}^{n}\frac{x_{k}-y_{k}}{x_{k+1}-y_{k+1}}\leqslant\sum_{k=1}^{n}\frac{1}{\lambda^{k-1}}=\frac{\lambda(1-\lambda^{-n})}{\lambda-1}<\frac{\lambda}{\lambda-1}.
\]
\end{enumerate}
\end{solution}

\begin{problem}
如图，以椭圆 \(\frac{x^2}{a^2}+\frac{y^2}{b^2}=1\)（\(a>b>0\)）的中心 \(O\) 为圆心，分别以 \(a\) 和 \(b\) 为半径作大圆和小圆. 过椭圆右焦点 \(F(c,0)\)（\(c>b\)）作垂直于 \(x\) 轴的直线交大圆于第一象限内的点 \(A\). 连结 \(OA\) 交小圆于点 \(B\). 设直线 \(BF\) 是小圆的切线.

\begin{enumerate}
\item 证明 \(c^2=ab\)，并求直线 \(BF\) 与 \(y\) 轴的交点 \(M\) 的坐标；
\item 设直线 \(BF\) 交椭圆于 \(P\)，\(Q\) 两点，证明：\(\overrightarrow{OP}\cdot\overrightarrow{OQ}=\frac{1}{2}b^2\).
\end{enumerate}

\bitmapfigure[max width=0.55\linewidth]{img/2006/tianjin_science/q22_fig1.png}
\end{problem}

\begin{solution}
椭圆的焦点满足 \(c^{2}=a^{2}-b^{2}\). 直线 \(x=c\) 与大圆 \(x^{2}+y^{2}=a^{2}\) 在第一象限的交点为
\[
A=(c,b),
\]
因为 \(c^{2}+b^{2}=a^{2}\). 又 \(OA=a\)，而 \(B\) 在射线 \(OA\) 上且在半径为 \(b\) 的小圆上，所以
\[
B=\frac{b}{a}A=\left(\frac{bc}{a},\frac{b^{2}}{a}\right).
\]
\begin{enumerate}
\item 直线 \(BF\) 是小圆在 \(B\) 处的切线，因此 \(OB\perp BF\)，即
\[
\overrightarrow{OB}\cdot\overrightarrow{FB}=0.
\]
等价地，\(B\cdot(F-B)=0\)，所以 \(B\cdot F=|B|^{2}=b^{2}\). 代入坐标得
\[
\frac{bc}{a}\cdot c=b^{2},
\]
从而 \(c^{2}=ab\).

小圆在 \(B(x_{B},y_{B})\) 处的切线方程为 \(x_{B}x+y_{B}y=b^{2}\)，代入 \(B\) 的坐标并约去 \(\dfrac{b}{a}\)，得
\[
cx+by=ab=c^{2}.
\]
令 \(x=0\)，得到直线与 \(y\) 轴的交点为 \(M=(0,a)\).
\item 令椭圆上 \(P\)，\(Q\) 的坐标分别为 \((x_{1},y_{1})\)，\((x_{2},y_{2})\)，并作变换 \(u=\frac{x}{a}\)，\(v=\frac{y}{b}\).
则 \(P\)，\(Q\) 对应于单位圆 \(u^{2}+v^{2}=1\) 上的两点，直线 \(BF\) 变为
\[
\alpha u+\beta v=1.
\]
其中 \(\alpha=\frac{c}{b}\)，\(\beta=\frac{b}{a}\).
令 \(t=\dfrac{b}{a}\)，由 \(c^{2}=ab\) 和 \(c^{2}=a^{2}-b^{2}\) 得 \(t^{2}+t-1=0\). 因而
\[
\alpha^{2}+\beta^{2}=\frac1t+t^{2}=2.
\]
直线 \(\alpha u+\beta v=1\) 到原点的距离为 \(1/\sqrt{2}\)，其垂足为 \((\alpha/2,\beta/2)\). 方向向量 \((-\beta,\alpha)\) 的长度为 \(\sqrt{2}\)，故单位圆上该弦的两个端点分别为 \(\left(\frac{\alpha+\beta}{2},\frac{\beta-\alpha}{2}\right)\) 和 \(\left(\frac{\alpha-\beta}{2},\frac{\beta+\alpha}{2}\right)\).
所以
\[
u_{1}u_{2}=\frac{\alpha^{2}-\beta^{2}}{4}.
\]
同理，\(v_{1}v_{2}=\frac{\beta^{2}-\alpha^{2}}{4}\).
于是
\[
\overrightarrow{OP}\cdot\overrightarrow{OQ}=x_{1}x_{2}+y_{1}y_{2}=(a^{2}-b^{2})\frac{\alpha^{2}-\beta^{2}}{4}.
\]
又 \(\alpha^{2}-\beta^{2}=\dfrac1t-t^{2}=2t\)，且 \(c^{2}=a^{2}t\)，故
\[
\overrightarrow{OP}\cdot\overrightarrow{OQ}=\frac{c^{2}t}{2}=\frac{a^{2}t^{2}}{2}=\frac{b^{2}}{2}.
\]
\end{enumerate}
\end{solution}
